\section{Perplexity Trade-offs, the Weaker-Links Account, and EWC}
\label{app:ppl-tradeoff}

\subsection{Perplexity trade-offs and the weaker-links hypothesis}

A recurring finding in the bilingual processing literature is that
proficiency in a speaker's two languages is not independent: gains in
one language are frequently accompanied by costs in the other. The
weaker-links account of \citet{gollan2008weaker} attributes this
pattern to frequency of use. Because a bilingual distributes a finite
amount of language experience across two systems, each word in each
language is encountered less often than it would be for a comparable
monolingual, so the links between lexical form and meaning are used
less and are correspondingly weaker. \citet{whitford2015second}
provide converging evidence from reading, showing that relative
exposure to each language modulates lexical processing speed in a
graded fashion, with the more-exposed language enjoying faster and more
reliable access.

Human-scale bilingual \beetle{} models reproduce the qualitative form
of this prediction. Under a fixed total experience budget, the amount
of effective exposure allocated to the second language (L2) is
determined by the curriculum, and increasing that allocation improves
L2 fluency while degrading first-language (L1) fluency. Because the two
languages draw on the same budget, encountering one language more often
necessarily means encountering the other less often, weakening the
under-exposed system in a manner directly analogous to the weaker-links
mechanism. Table~\ref{tab:ppl-curriculum} reports validation perplexity
by curriculum for two language pairs, Dutch--English and
Chinese--English, together with monolingual controls trained on the
same budget.

\begin{table}[t]
\centering
\small
\begin{adjustbox}{max width=\columnwidth}
\begin{tabular}{@{}lcccc@{}}
\toprule
& \multicolumn{2}{c}{Dutch model} & \multicolumn{2}{c}{Chinese model} \\
\cmidrule(lr){2-3}\cmidrule(lr){4-5}
Curriculum & \makecell{Dutch\\PPL (L1)} & \makecell{English\\PPL (L2)} & \makecell{Chinese\\PPL (L1)} & \makecell{English\\PPL (L2)} \\
\midrule
B1 balanced      & 318.4 & 412.4  & 5198.1 & 395.4  \\
B2 simultaneous  & 273.8 & 1128.8 & 7550.8 & 1309.1 \\
B3 sequential    & 268.4 & 1032.8 & 7222.9 & 1277.4 \\
B4 classroom     & 274.8 & 805.6  & 7073.0 & 948.8  \\
B5 late          & 268.3 & 987.4  & 7729.2 & 1252.3 \\
\midrule
Monolingual      & 323.7 & 781.2  & 3839.4 & 781.2  \\
\bottomrule
\end{tabular}
\end{adjustbox}
\caption{Validation perplexity by curriculum for Dutch--English and
Chinese--English \beetle{} models, with monolingual controls trained on
a matched budget. For each language pair the L1 and L2 columns move in
opposite directions as effective L2 exposure changes for the
typologically close pair, but move together for the distant pair.}
\label{tab:ppl-curriculum}
\end{table}

For Dutch--English the results trace out the expected trade-off. The
balanced curriculum B1, which carries the highest effective English
exposure, attains an English perplexity of roughly $412$, about
$2.5\times$ lower than the staged curricula, which cluster near
$1050$. The same curriculum pays for this with a Dutch perplexity of
$318$ against roughly $270$ for the staged alternatives, an increase of
about $18\%$. The two languages therefore move in opposite directions
along a single shared experience axis: shifting the budget toward
English improves English at the expense of Dutch, exactly the pattern
the weaker-links account predicts for two closely related languages
that share substantial lexical and orthographic structure.

Chinese--English falls in a different regime. Here every bilingual
curriculum is worse than the corresponding monolingual control in
\emph{both} languages: Chinese perplexity rises to roughly $7500$
against $3839$ for the monolingual control, and English perplexity sits
near $1280$ against $781$. Rather than trading L1 fluency for L2
fluency along a shared axis, splitting the budget between two systems
that share little lexical or orthographic structure incurs a dual cost,
degrading both languages relative to what the same budget buys when
spent on either alone. The trade-off characteristic of the weaker-links
prediction thus appears to depend on the two languages sharing enough
structure to compete over a common representational substrate; when
that overlap is absent, division of the budget is simply lossy on both
sides. These comparisons are intended as a quick quantitative check
against the weaker-links prediction rather than a full account of the
underlying representations.

\subsection{Comparison with Constantinescu et al.\ (2025): does EWC prevent L1 attrition?}

A natural question is whether an explicit regularisation constraint can
protect the first language during sequential L1$\rightarrow$L2 training.
\citet{constantinescu2025investigating} examine this question directly
and report that critical-period-like retention does not emerge on its
own: sequential exposure drives severe attrition of the first language
unless training is constrained to preserve it. The \beetle{} models are
consistent with this conclusion. Under the B3 sequential curriculum,
first training on L1 and then on L2, L1 perplexity degrades
catastrophically, and adding an elastic weight consolidation (EWC)
penalty \citep{kirkpatrick2017overcoming} substantially mitigates that
degradation. The effect, however, is weaker at the $24$B-token scale
studied here than the near-complete retention reported by
\citet{constantinescu2025investigating}.

Table~\ref{tab:ewc-24b} contrasts the sequential curriculum with and
without EWC at $24$B tokens for Chinese--English, using a penalty
strength of $\lambda = 20$. EWC lowers final Chinese (L1) perplexity by
roughly two orders of magnitude while raising English (L2) perplexity
only slightly, from $86.9$ to $94.7$.

\begin{table}[t]
\centering
\small
\begin{adjustbox}{max width=\columnwidth}
\begin{tabular}{@{}lrrl@{}}
\toprule
Model & \makecell[r]{L1 (Chinese)\\PPL} & \makecell[r]{L2 (English)\\PPL} & Interpretation \\
\midrule
EWC ($\lambda{=}20$) & $437{,}680$    & $94.7$ & \makecell[l]{L1 forgetting reduced,\\small L2 cost} \\
\addlinespace
Non-EWC (B3) & $14{,}295{,}664$ & $86.9$ & \makecell[l]{Catastrophic L1\\forgetting} \\
\bottomrule
\end{tabular}
\end{adjustbox}
\caption{$24$B Chinese--English sequential training with and without
EWC ($\lambda = 20$). EWC reduces L1 forgetting by about two orders of
magnitude at a small L2 penalty.}
\label{tab:ewc-24b}
\end{table}

Because raw perplexities in the millions are difficult to interpret,
Table~\ref{tab:ewc-nats} restates the same comparison in nats, giving
the best L1 loss reached during training, the final L1 loss, the
resulting forgetting (the increase from best to final), and the
implied final L1 perplexity.

\begin{table}[t]
\centering
\small
\begin{adjustbox}{max width=\columnwidth}
\begin{tabular}{@{}lrrrr@{}}
\toprule
Model & \makecell[r]{Best L1\\(nats)} & \makecell[r]{Final L1\\(nats)} & \makecell[r]{For-\\getting} & \makecell[r]{Final L1\\PPL} \\
\midrule
Non-EWC & $7.43$ & $14.08$ & $+6.66$ & $18{,}957{,}016$ \\
EWC     & $7.30$ & $13.50$ & $+6.20$ & $4{,}697{,}590$ \\
\bottomrule
\end{tabular}
\end{adjustbox}
\caption{L1 forgetting measured in nats. EWC reduces final Chinese
forgetting by about $0.46$ nats, corresponding to a roughly $4\times$
lower final L1 perplexity, while the best L1 loss reached is nearly
unchanged.}
\label{tab:ewc-nats}
\end{table}

Table~\ref{tab:ewc-retention} tracks L1 perplexity across the L2 phase
of training, reporting the EWC and non-EWC values at matched steps
together with their ratio. The ratio falls well below one throughout,
reaching a minimum near $0.04\times$ mid-phase before the two
trajectories partially reconverge, leaving EWC at roughly a quarter of
the non-EWC L1 perplexity at the final step.

\begin{table}[t]
\centering
\small
\begin{adjustbox}{max width=\columnwidth}
\begin{tabular}{@{}lrrr@{}}
\toprule
Step & Non-EWC & EWC & \makecell[r]{EWC /\\Non-EWC} \\
\midrule
$32768$ (pre)    & $3{,}574$       & $2{,}727$       & $0.76\times$ \\
$45775$ (phase)  & $175{,}945$     & $68{,}305$      & $0.39\times$ \\
$47824$          & $599{,}079$     & $159{,}117$     & $0.27\times$ \\
$49872$          & $1{,}973{,}516$ & $259{,}442$     & $0.13\times$ \\
$53968$          & $7{,}906{,}358$ & $1{,}107{,}082$ & $0.14\times$ \\
$62160$          & $51{,}264{,}313$& $2{,}089{,}939$ & $0.04\times$ \\
$78544$ (final)  & $18{,}957{,}016$& $4{,}697{,}590$ & $0.25\times$ \\
\bottomrule
\end{tabular}
\end{adjustbox}
\caption{L1 (Chinese) perplexity across the L2 training phase, with and
without EWC, and their ratio. EWC consistently retains more of the
first language, but does not restore it to its pre-phase level.}
\label{tab:ewc-retention}
\end{table}

Taken together, these measurements show that EWC reduces final Chinese
forgetting by about $0.46$ nats, equivalent to a roughly $4\times$
lower matched-step L1 perplexity, at a small English penalty of about
$0.10$ nats. The constraint therefore does what its motivation predicts:
it preserves parameters important to the first language and slows
attrition. Yet Chinese remains substantially degraded relative to its
best value, in contrast to the near-initial retention reported by
\citet{constantinescu2025investigating}. A likely explanation is that
$\lambda = 20$, a penalty strength calibrated on smaller models, is
under-regularised at $24$B scale: as model and data scale grow, the
Fisher-weighted penalty required to hold important parameters in place
grows with them, so a fixed $\lambda$ exerts proportionally less
pressure. Establishing the penalty strength needed to reproduce strong
retention at scale would require a scale-aware sweep over $\lambda$
rather than transfer of a value tuned on smaller runs. The broader
conclusion is that preventing catastrophic forgetting at scale is not
simply a matter of switching on an EWC term; it requires scale-aware
calibration of the constraint itself.
